\bookchapter{פרק 24: ניתוח מעמיק של הרשאות Root ונעילת BL}

\chapterintro{פרק זה הוא הרחבה מעמיקה של תוכן פרק 20, מתאים למשתמשים שרוצים באמת להבין את מנגנוני השכבה הנמוכה של הטלפון, ואפילו לנסות להתקין ROM. שים לב, פעולות אלה כרוכות בסיכון מסוים, אנא נהג בזהירות.}

\booksection{הבנה מעמיקה של Bootloader}

אימות חתימה

\booksubsection{: נעילת BL מאמתת את החתימה הדיגיטלית של תמונת האתחול, רק מערכת חתומה על ידי היצרן יכולה לעלות.}

\begin{itemize}
  \item \textbf{מניעת ירידת גרסה}: מונעת מהמשתמש להתקין גרסת מערכת ישנה כדי לנצל פרצות ידועות.
  \item \textbf{הגנה על נתונים}: גם אם מותקנת מערכת חדשה, ללא סיסמת החשבון המקורית לא ניתן לגשת לנתוני המשתמש.
  \item \textbf{2. תהליך מפורט לפתיחת נעילת BL (דוגמת Xiaomi)}בקשת הרשאת פתיחה
\end{itemize}

\booksubsection{אחסון ואימות ברמת חומרה}

BL (Locked/Unlocked) ומפתח יצרן, בחומרה מוגנת (devinfo/eFuse/RPMB). בכל הפעלה, Bootloader קורא ומחליט אם לבצע אימות חתימה.

\booksubsection{אימות חתימה דיגיטלית}

יצרן חותם במפתח פרטי. בהפעלה, Bootloader מאמת עם מפתח ציבורי (boot.img, recovery.img). אופן:

\begin{enumerate}
  \item חישוב מחדש של ערך hash לתמונת מערכת.
  \item פענוח חתימה עם מפתח ציבורי = hash.
  \item השוואת hash, זהה = עובר. מבטיח שלמות.
\end{enumerate}

\booksubsection{שרשרת אתחול מאובטח}

BL = החוליה הראשונה. אימות טעינת kernel שרשרת אמון.

\booksubsection{נעילה מול פתיחה: ההשפעה הדו-כיוונית של נעילת BL}

להבנת BL, טבלת השוואה ב-5 ממדים:

\begin{center}
\begin{tabularx}{\linewidth}{|X|X|X|}
\hline
\textbf{תפקיד מרכזי} & \textbf{מצב נעול (ברירת מחדל)} & \textbf{מצב לאחר פתיחה} \tabularnewline
\hline
אבטחת מערכת & מניעת צריבת ROM לא רשמי, הגנת שלמות\newline בטוח: רק רשמי, נגד חבלה & פחות בטוח: ROM צד ג', אימות מבוטל, Root = חבלה \tabularnewline
יציבות מערכת & מערכת רשמי, יציב\newline יציב: רשמי, OTA & סיכון: ROM צד ג', קריסה/brick, חומרה, OTA נכשל \tabularnewline
זכויות אחריות & BL נעול = אחריות\newline אחריות: רשמי & אחריות בוטלה: BL = ויתור, גם נעילה מחדש ניתן לזהות \tabularnewline
אבטחת נתונים & פתיחה = מחיקת נתונים\newline בטוח: BL מגביל חבלה & סיכון: מחיקה, הצפנה, איבוד = חילוץ \tabularnewline
מגבלות פונקציה & תשלום/בנק עובד\newline פונקציונליות מלאה: כל הפונקציות הרשמיות פעילות & מוגבל: Google Pay/בנק דוחים, DRM יורד \tabularnewline
\hline
\end{tabularx}
\end{center}

\booksection{תהליך מפורט לפתיחת נעילת BL (דוגמת Xiaomi)}

\begin{enumerate}
  \item \textbf{קישור חשבון}: בהגדרות הטלפון קשר חשבון Xiaomi.
  \item \textbf{תקופת המתנה}: בהתאם לרמת החשבון, יש להמתין 7-30 ימים.
  \item \textbf{פעולת פתיחה}: השתמש בכלי פתיחת הנעילה של Xiaomi בחיבור למחשב לפתיחת הנעילה.
  \item \textbf{שלוש. אפשרויות Root שונות}Magisk
\end{enumerate}

\booksection{שלוש. אפשרויות Root שונות}

\begin{itemize}
  \item \textbf{Magisk}: פתרון Root מבוסס-ליבה, עמוק יותר, הסתרה טובה יותר.
  \item \textbf{SuperSU}: פתרון Root ותיק, הפסיק להתעדכן.
  \item \textbf{ארבע. מבוא להתקנת ROM}הבהוב (התקנת ROM) פירושו התקנת מערכת הפעלה חדשה בטלפון. סוגי הבהוב נפוצים:
\end{itemize}

\booksection{ארבע. מבוא להתקנת ROM}

: התקנת מערכת רשמית, להצלת מכשיר או הורדת גרסה.

\begin{itemize}
  \item \textbf{הבהוב Recovery צד-שלישי}: כמו TWRP, לגיבוי ושחזור, הבהוב ROM צד-שלישי.
  \item \textbf{הבהוב ROM צד-שלישי}: כמו LineageOS, Pixel Experience, לחוות מערכות שונות.
  \item \textbf{אזהרה חשובה}הבהוב כרוך בסיכון, יש לפעול בזהירות. עלול לגרום ל:
\end{itemize}

\booksection{ה. Samsung Knox: המחיר הבלתי-הפיך של פתיחת BL}

Samsung: “נתיק” ייחודי:\textbf{Samsung: Knox נשרף, Samsung Pay אובד, נתונים נמחקים, אחריות.}. בפרט:

\begin{itemize}
  \item \textbf{הנתיך חד-פעמי}:בתוך הטלפון יש נתיך eFuse, ברגע שמזוהה ש-BL שוחרר הנתיך נשרף לצמיתות, לא ניתן לשחזר.
  \item \textbf{פונקציות מושפעות}:Samsung Pay (כרטיס אוטובוס/בנק), תיקיית Knox מאובטחת, ניהול MDM ארגוני, אפליקציות ביטוח בריאות - כולם יפסיקו לעבוד.
  \item \textbf{מחיקת נתונים בכפייה}:לאחר הפעלת הנתיק הטלפון יבצע איפוס יצרן, כל הנתונים המקומיים יימחקו (נתוני ענן לא מושפעים).
  \item \textbf{השפעה על אחריות}:Huawei מצהירה רשמית שלאחר נתיק Knox אין אחריות.
\end{itemize}

\begin{tcbwarn}{חשבו פעמיים לפני שתפעלו}
הנתיך של Samsung Knox הוא\textbf{בלתי הפיך}לאחר ההפעלה לא ניתן לחזור למצב יצרן. אם רק רצית פונקציות הדורשות Root, מומלץ לבדוק חלופות ללא Root.

\end{tcbwarn}

\textbf{הטלפון הופך ל“לבנה” (לא נדלק)}\par

אובדן אחריות

\begin{itemize}
  \item אובדן נתונים
  \item תקלה בפונקציות תשלום
  \item לפני הבהוב יש לבצע גיבוי מלא, לקרוא בעיון את המדריך, ולוודא התאמה לדגם ולגרסה.
  \item משאבים מומלצים
\end{itemize}

XDA Developers

\begin{tcbinfo}{: קהילת מפתחי האנדרואיד הגדולה בעולם, עם שפע מדריכי הבהוב ומשאבי ROM.}
\begin{itemize}
  \item \textbf{Coolapk (CoolAPK)}: קהילה דיגיטלית ידועה, עם מדריכים ודיונים רבים.
  \item \textbf{פורומים לפי דגם}: כמו קהילת Xiaomi, קהילת OnePlus, עם מדריכים רשמיים וקהילתיים.
  \item \textbf{פרק 25: טיפול בטלפון יד שנייה ומורשת דיגיטלית}היי, חברים, אחרי יותר מעשרים פרקים של "אימון", אני מאמין שכבר הפכתם מ"טירון דיגיטלי" למשתמש טלפון די בטוח בעצמו. אבל לחיי הטכנולוגיה יש תמיד מחזוריות – החדש הופך בסוף לישן. כשאתם מתכוננים "להפריש" את החבר הוותיק, האם אתם יודעים איך להיפרד בכבוד ובבטחה?
\end{itemize}

\end{tcbinfo}

\booksection{ו. שאלות נפוצות (Q\& A)}

\textbf{Q1: האם ניתן לשחזר לאחר פתיחת נעילת BL?}\par

\textbf{A:}ניתן לנעול מחדש. בדרך כלל נדרש לצרוב חבילת יצרן מלאה ולהשתמש בסקריפט צריבה (כמו של Xiaomi)\texttt{flash\_all\_lock.\allowbreak bat}) כדי לנעול. אך שים לב, נעילה מחדש לא מבטלת נתיכי אבטחה (כמו Knox של Samsung).

\textbf{Q2: איך מבצעים Root לאחר פתיחה?}\par

\textbf{A:}פתיחת BL היא רק הצעד הראשון. Root בדרך כלל דורש:

\begin{enumerate}
  \item צרוב Recovery מותאם (כמו TWRP).
  \item דרך Recovery צרוב Magisk לתיקון תמונת boot והשגת Root.
\end{enumerate}
